\documentclass[11pt]{article}
\usepackage[a4paper,margin=27mm]{geometry}
\usepackage{amsmath,amssymb,mathtools}
\usepackage{fontspec,unicode-math}
\setmainfont{Latin Modern Roman}
\setmathfont{Latin Modern Math}
\usepackage{xurl,hyperref}
\hypersetup{colorlinks=true,urlcolor=blue,linkcolor=blue}
\newcommand{\sref}[2]{\hyperlink{source:#1:#2}{[S#2]}}
\newcommand{\cataloguescope}[1]{\paragraph{Mathematical question.} #1}
\begin{document}
% UnsolvedMath ID 8400021; source catalogue code AMR-083-0021
\setcounter{section}{8400020}
\section{O18b — Recognizing primitive roots randomly}\label{8400021}
{\small\sffamily UnsolvedMath ID 8400021 · Graph Theory}
\subsection{Definitions and mathematical statement}

\paragraph{Shared notation.}$\mathbb N=\{1,2,\ldots\}$, and $\mathbb Z,\mathbb Q,\mathbb R,\mathbb C$ denote the integers, rationals, reals and complex numbers. Any other conventions are specified locally.


Integers are supplied in binary; polynomial time means a polynomial number of classical bit operations in total input length. On promised inputs $b\in\mathbb N$ and a prime $p$, call $b$ a primitive root modulo $p$ if its residue generates $(\mathbb Z/p\mathbb Z)^\times$. The source's notion of recognition in $R$ means one-sided randomized acceptance: positive instances are accepted with probability at least $1/2$, negative instances are never accepted, and other outcomes are designated $?$.
\paragraph{Statement recorded in the supplied source.}
Can primitive roots modulo a prime be recognized in $R$ in this sense, without receiving the factorization of $p-1$? \sref{8400021}{2}
\subsection{Short English statement}
Is there a polynomial-time randomized procedure which can certify primitive-root instances with bounded success probability and never falsely certify a non-generator?
\subsection{Sources}
\begin{description}
\item[\hypertarget{source:8400021:1}{S1}] ULAM.AI and the UnsolvedMath Contributors, UnsolvedMath: A Curated Collection of Open Mathematics Problems, record 8400021 (AMR-083-0021). CC BY 4.0. \url{https://huggingface.co/datasets/ulamai/UnsolvedMath}.
\item[\hypertarget{source:8400021:2}{S2}] Leonard M. Adleman and Kevin S. McCurley, Open Problems in Number Theory Complexity, II, in Algorithmic Number Theory, LNCS 877 (1994), pp. 291--322, C18/O18b, manuscript p. 14 (LNCS p. 304), and Definition 5. \url{https://mccurley.org/papers/open.ps.gz}.
\end{description}
\label{end:8400021}
\end{document}
